\chapter{Fazit und Ausblick}
\label{cha:Fazit und Ausblick}

\section{Beantwortung der Forschungsfragen}

Die vorliegende Arbeit untersuchte, inwiefern ein RL-basierter Ansatz eine
autonome Drohne befähigen kann, mithilfe lokaler Tiefensensorik in
hindernisreichen Umgebungen sicher zu landen.

Der trainierte Agent erlernt eine reaktive Ausweichstrategie, deren
Leistung primär durch strukturelle Entwurfsentscheidungen begrenzt wird:
Sensorhorizont, Entscheidungsintervall und Belohnungsgewichtung bestimmen,
wie knapp der Agent Hindernisse passiert. Die Wahl der Kollisionsmetrik
erweist sich als entscheidender Faktor: Unter oberflächenbasierter
Erkennung erreicht der Agent 98,4\,\% und übertrifft damit den
privilegierten Pfadplaner (82,0\,\%).

Die Transferuntersuchung zeigt eine klare Trennung zweier
Teilfähigkeiten: Hindernisnavigation generalisiert auf die
Weinbergumgebung (69,7\,\% bei erweitertem Erfolgsradius); die terminale
Landephase scheitert an der im Training nie erfahrenen Konfiguration von
Vegetation in der Landezone, verstärkt durch Perceptual Aliasing und
fehlendes zeitliches Gedächtnis.

Der Vergleich mit dem privilegierten Pfadplaner quantifiziert den Vorteil
vollständiger Umgebungskenntnis. Unter fairer Kollisionsmetrik löst der
lernbasierte Ansatz das Navigationsproblem im Korridor erfolgreicher als
der Planer, bei um Größenordnungen niedrigerer Inferenzzeit. Im Weinberg
kehrt sich dieses Verhältnis um; die Leistungslücke reflektiert die
begrenzte Trainingsdiversität.

Insgesamt zeigt die Arbeit, dass ein RL-Agent mit minimaler Sensorik die
kombinierte Aufgabe aus Zielnavigation und Hindernisvermeidung in
Simulation lösen kann und die erlernte Strategie sich teilweise auf ein
realistisches Szenario übertragen lässt.

\section{Perspektiven für zukünftige Forschung}

Eine Multi-Seed-Validierung und eine Skalierung auf größere
Trainingsbatches könnten die statistische Belastbarkeit und die
Strategiequalität verbessern. Eine rekurrente Architektur oder ein
Tiefenbild-Stapel würde dem Agenten zeitliches Gedächtnis verleihen und
das Perceptual-Aliasing-Problem adressieren. Eine Ablationsstudie könnte
den individuellen Beitrag von Tiefenbild und Zustandsvektor
quantifizieren. Domain Randomization, insbesondere Sensorrauschen und
variable Zielpositionen, könnte die Robustheit und den Transfer auf reale
Hardware vorbereiten. Langfristig stellt der Sim-to-Real-Transfer den
entscheidenden nächsten Schritt dar, um die gezeigte Leistungsfähigkeit
in realen Einsatzszenarien zu validieren.
